\chapter{Conclusion}
\label{ch:conclusion}

TalkCS was developed as a BSc thesis project at the Faculty of Informatics of Eötvös
Loránd University, with the goal of giving university students a single, structured
place for academic collaboration instead of scattered social media groups and chat
channels. The resulting platform combines a question-and-answer forum organized by
subjects, threaded and live-updating discussions, accepted answers, a voting and
reputation system with levels, badges and a leaderboard, a per-subject resource library,
global search and trending content, user profiles with following, mentions and
notifications, real-time direct messaging, a shared academic calendar with a review
workflow, and an administration panel for analytics, role management and configuration.

The system was implemented as a three-tier application with a Spring Boot backend, a
PostgreSQL database and a React single-page frontend. The backend exposes a documented
REST API secured with stateless JWT authentication and role-based authorization, and a
STOMP-over-WebSocket channel for real-time features. To prevent duplicate questions, a
similarity detection algorithm based on the BM25 ranking function was designed and
implemented, which combines the textual relevance of the post body with the similarity
of titles and tags and runs while the user is still typing. The correctness of the
backend is supported by 205 automated tests covering services, controllers, repositories
and the full authentication flow, with 69\% instruction coverage, and the frontend is
covered by 100 tests of its API clients, pages and hooks.

\section*{Future Work}

While the core functionality is complete, several directions could further improve the
platform:

\begin{itemize}
  \item \textbf{Semantic similarity:} the lexical BM25 approach recognizes duplicates that
    use the same terminology. Sentence embeddings produced by a language model would also
    recognize questions that are phrased with entirely different words, and could be
    used for smarter tag suggestions as well.
  \item \textbf{Richer resource library:} written reviews of resources in addition to
    votes, and organizing materials by semester or academic year.
  \item \textbf{Moderation tools:} reporting of inappropriate content and a moderation
    queue for professors.
  \item \textbf{Configurable notifications:} letting users choose which notification types
    they want to receive by e-mail, and watching individual posts.
  \item \textbf{Branding in the interface:} applying the stored site name, tagline, colour
    and logo throughout the user interface, so that each institution's installation
    carries its own identity.
  \item \textbf{Production hardening:} versioned database migrations (e.g. with Flyway)
    instead of automatic schema updates, object storage for uploaded files, PostgreSQL
    full-text search for large forums, and rate limiting of the
    authentication endpoints.
  \item \textbf{Mobile application:} because all functionality is available through the
    REST and WebSocket API, a native mobile client could be built on the existing backend
    without server-side changes.
\end{itemize}

Building TalkCS provided extensive hands-on experience with the full life cycle of a web
application: designing a relational data model, building a secure layered backend with
Spring, implementing real-time communication, developing a modern TypeScript frontend,
applying a classical information retrieval algorithm to a practical problem, and testing
the system at several levels. The project demonstrates that a feature-rich and practical
academic platform can be delivered by a single developer within the timeframe of a
thesis, and it provides a solid foundation that a university could deploy and extend.
